\section{讨论与开放项：哪些量子结构需要额外输入？}

\subsection{本文已闭合的最小结论}

本文的核心结论是：\textbf{粗结果概率由纤维计数给出}，而 Born 权重为该概率在有效 Hilbert 表述下的等价写法（定理 \ref{thm:born_as_counting_equivalence}）。
此外，POVM 与去相干通道被用作有限读出的标准接口语言，用以刻画更一般的读出与粗粒化操作。

\subsection{必须显式增加的输入（否则无法推进）}

以下结构并非由折叠纤维计数自动推出；若要进一步获得更强的量子结构陈述，需显式增加相应的最小桥接假设或将其列为开放项：
\begin{itemize}
  \item \textbf{复线性与复数幅度：} 本文采用 $\sqrt{p_m(x)}e^{i\theta(x)}$ 的复幅度参数化；复 Hilbert（相对于实/四元数情形）的选择需额外原则（例如对称性与复合结构）。
  \item \textbf{复合系统与张量积：} 折叠本体在本稿中以单一窗口读出为中心；张量积结构需要关于“并置/独立读出通道”的协议性输入。
  \item \textbf{动力学与幺正演化：} 本文处理静态读出概率；幺正/半群形式的动力学陈述需要额外的可逆性、生成元、或可审计更新规则。
  \item \textbf{Born 形式的刚性刻画：} 若要求 Born 形式为概率赋值的唯一可能形式，可引入投影/效应上的非语境有限可加性并使用 Gleason--Busch 类型定理；其假设强度与适用边界需另行审计。
\end{itemize}

\subsection{与离散连接/holonomy 的潜在拼接点}

折叠本体中已有离散连接与 holonomy 的语义位置：沿回路的平行移动复合可产生不可消去的历史依赖。
若未来要把“相位”从任意函数 $\theta(x)$ 升级为可检验结构，一个自然方向是：
\begin{itemize}
  \item 选择具体的残差载体（例如 $U(1)$ 或其有限近似）作为纤维群；
  \item 将平行移动实现为该群上的作用，从而使 holonomy 直接给出相位回路不变量；
  \item 设计可干涉的读出协议，使得非对角效应可被稳定估计，从而让相位不再是“规约自由度”。
\end{itemize}

